% Generated draft: requires mathematical and source review.
\documentclass[11pt]{article}
\usepackage[margin=25mm]{geometry}
\usepackage{fontspec}
\setmainfont{Latin Modern Roman}
\usepackage{amsmath,amssymb,mathtools,unicode-math}
\setmathfont{Latin Modern Math}
\usepackage{enumitem,needspace,xurl,hyperref}
\setlength{\parindent}{0pt}
\setlength{\parskip}{5pt}
\setlength{\emergencystretch}{3em}
\newcommand{\N}{\mathbb{N}}
\newcommand{\Z}{\mathbb{Z}}
\newcommand{\Q}{\mathbb{Q}}
\newcommand{\R}{\mathbb{R}}
\newcommand{\C}{\mathbb{C}}
\newcommand{\F}{\mathbb{F}}
\newcommand{\A}{\mathbb{A}}
\newcommand{\PP}{\mathbb{P}}
\newcommand{\E}{\mathbb{E}}
\newcommand{\eps}{\varepsilon}
\newcommand{\dd}{\,\mathrm{d}}
\begin{document}
\section*{UnsolvedMath ID: 11000119}
\textbf{Source problem code:} catalogue entry
\subsection{Definitions and mathematical statement}
Let \(N_h=\#^h\R P^2\) be a closed connected nonorientable surface of genus \(h\), and let \(\operatorname{Mod}(N_h)=\pi_0\operatorname{Diff}(N_h)\) be its full mapping class group. No orientation-preserving subgroup is selected on this nonorientable surface. For the genera intended in the source, the compound problem asks three generation questions:
\begin{enumerate}
\item Do there exist \(a,b\in\operatorname{Mod}(N_h)\) with \(\langle a,b\rangle=\operatorname{Mod}(N_h)\)?
\item Can \(a,b\) be chosen to have finite order?
\item Do there exist involutions \(u,v,w\), each of order two, with \(\langle u,v,w\rangle=\operatorname{Mod}(N_h)\)?
\end{enumerate}
The second is a two-torsion-generator question, not merely a finite generating set consisting of torsion elements. The third fixes the number three and requires order-two elements; it does not restrict them to any particular geometric kind of involution. These conditions are on generators of the entire mapping class group, not a finite-index subgroup.

The source does not state a genus threshold or exceptions, nor whether ``by two/three'' means at most that many generators or a set of that exact size. The displayed existence statements use the customary bounded-generator reading without demanding that the group have minimal rank exactly two or three. Small-genus cases and whether identity is allowed in a torsion tuple need the intended convention; genuine involutions above have order exactly two. Thus no single all-\(h\geq1\) affirmative assertion is silently appended.
\subsection{Short English statement}
Can a closed nonorientable surface's full mapping class group be generated by two elements, by two finite-order elements, or by three involutions? The genus range and small-case conventions are not supplied.
\subsection{Sources}
\begin{itemize}
\item Local UnsolvedMath registry, permanent ID \(11000119\), exact source code \texttt{catalogue entry}.
\item Farb, \emph{Problems on Mapping Class Groups and Related Topics}, Problem 3.1, PDF page 98, parts (a)--(b): \url{https://www.math.uchicago.edu/~farb/papers/mcgbook.pdf}. Registry access 2026-07-29; pdf-font-structure extraction.
\item The research summary refers generally to later generation work without any precise citation.
\end{itemize}
\end{document}
